\documentclass[conference]{IEEEtran}
\IEEEoverridecommandlockouts

\usepackage{cite}
\usepackage{amsmath,amssymb,amsfonts}
\usepackage{graphicx}
\usepackage{textcomp}
\usepackage{xcolor}
\usepackage{booktabs}
\usepackage{array}
\usepackage{url}
\usepackage{algorithm}
\usepackage{algorithmic}
\usepackage{float}
\usepackage{multirow}

\begin{document}

\title{Designing an India-Centric Operational Healthcare Database for Machine Learning and Predictive Donor Response Modeling in BloodLink AI}

\author{
    \IEEEauthorblockN{Lucky Dadlani}
    \IEEEauthorblockA{\textit{Department of Information Science and Engineering} \\
    \textit{CMR Institute of Technology}\\
    Bengaluru, India \\
    Email: luckydadlani@example.com}
    \and
    \IEEEauthorblockN{Nivedita Vinod}
    \IEEEauthorblockA{\textit{Department of Information Science and Engineering} \\
    \textit{CMR Institute of Technology}\\
    Bengaluru, India \\
    Email: niveditavinod@example.com}
}

\maketitle

\begin{abstract}
Predictive machine learning in healthcare operations critically depends on the architectural fidelity and semantic structure of the underlying database. In emergency blood donation systems, operational databases must concurrently support low-latency transactional workflows and serve as high-quality analytic substrates for predictive modeling. This paper presents the database architecture, realistic dataset generation methodology, and predictive modeling pipeline of BloodLink AI---an integrated emergency blood coordination platform connecting hospitals, blood banks, and donors across Indian urban centers. We propose a normalized, India-centric relational schema implemented on Supabase PostgreSQL that captures spatial coordinates, component-level inventory, emergency lifecycles, and granular notification logs. Using 426 live operational records, we establish a rigorous machine learning benchmark comparing six classification algorithms (Logistic Regression, Support Vector Machines with RBF kernel, $k$-Nearest Neighbors, Decision Trees, Gradient Tree Boosting, and Random Forest) under both an 80/20 stratified split and 10-fold cross-validation. Random Forest demonstrates optimal operational stability (Accuracy $0.8053 \pm 0.0518$, ROC-AUC $0.8366 \pm 0.0899$, $F_1$-score $0.8634 \pm 0.0372$), markedly outperforming heuristic distance-only baselines (ROC-AUC $0.5101$). Building on these empirical results, we formulate a generalized algorithm selection framework for healthcare tabular problems categorized by data dimensionality, class imbalance, and latency constraints. The study demonstrates that aligning operational database engineering with machine learning requirements is essential for reliable clinical decision support.
\end{abstract}

\begin{IEEEkeywords}
Healthcare informatics, emergency blood supply chain, operational database design, predictive donor modeling, Random Forest, tabular machine learning, Supabase PostgreSQL.
\end{IEEEkeywords}

% =========================================================================
\section{Introduction}
% =========================================================================
Emergency blood coordination is fundamentally an information and logistical optimization problem before it becomes a machine learning task \cite{belien2012supply}. When a medical emergency arises in a hospital, clinical staff must rapidly resolve three interdependent operational queries: (1) which regional blood banks possess viable, component-compatible inventory; (2) which eligible voluntary donors are geographically proximate to the requesting facility; and (3) which among those donors possesses the highest empirical probability of responding positively and promptly to an urgent dispatch \cite{shringi2020blood}. 

In resource-constrained and densely populated healthcare ecosystems such as India, traditional blood coordination relies predominantly on fragmented registries, unindexed emergency broadcast lists, and manual telephonic outreach \cite{agarwal2012blood, gemawat2021blood}. This manual approach induces substantial operational friction: emergency response times frequently extend into hours, non-targeted donor broadcasting causes acute communication fatigue, and concurrent donor confirmations lead to unmanaged reservation conflicts \cite{laermans2022use}. 

Modern machine learning (ML) models offer significant potential for prioritizing donor outreach by predicting donor responsiveness from historical, spatial, and temporal features \cite{kauten2022predicting, yeh2024prediction}. However, a critical disconnect often persists in medical informatics between operational database engineering and downstream predictive modeling. Conventional systems frequently store transactional records in formats unsuitable for feature extraction or train algorithms on synthetic, ungrounded benchmark datasets that fail to reflect real-world clinical distributions, spatial constraints, and timing irregularities \cite{walonoski2018synthea, choi2017medgan}.

To address these challenges, this paper presents the complete architectural and predictive modeling framework of \textbf{BloodLink AI}, a full-stack emergency blood coordination platform. The core contributions of this work are twofold:
\begin{enumerate}
    \item \textbf{An Operational Medical Database Architecture for ML:} We introduce a normalized, India-centric relational schema implemented on Supabase PostgreSQL that unifies identity, multi-component inventory, emergency lifecycle tracking, and communication telemetry. We formalize principles for realistic medical dataset generation, detailing how spatial heterogeneity, component-level inventory distributions, and non-random missingness patterns are preserved.
    \item \textbf{A Generalized Healthcare Algorithm Selection Framework:} We provide a comprehensive taxonomy guiding algorithm selection across distinct healthcare tabular data scenarios (binary outcome classification, survival/time-to-event modeling, severe class imbalance, cold-start deployments, and high-dimensional spaces). We validate this framework through a rigorous multi-model empirical benchmark on 426 live operational records, evaluating six classifiers across single-split and 10-fold cross-validation protocols against non-ML heuristic baselines.
\end{enumerate}

The remainder of this paper is structured as follows: Section II surveys related work in medical data synthesis, blood supply chains, and tabular healthcare ML. Section III details the system and database architecture. Section IV presents our dataset generation and noise modeling principles. Section V describes the feature engineering pipeline. Section VI introduces the generalized algorithm selection framework. Section VII details empirical experimental results. Sections VIII, IX, and X discuss operational implications, limitations, and ethical considerations, and Section XI concludes the paper.

% =========================================================================
\section{Related Work}
% =========================================================================

\subsection{Synthetic and Realistic Medical Dataset Generation}
Generating synthetic healthcare records that accurately mimic real clinical environments while protecting patient privacy is an active research domain. Walonoski et al. \cite{walonoski2018synthea} developed \textit{Synthea}, an open-source synthetic patient generator that simulates realistic clinical lifecycles without disclosing Protected Health Information (PHI). In the deep learning literature, Choi et al. \cite{choi2017medgan} proposed \textit{medGAN}, using generative adversarial networks with autoencoders to generate realistic discrete electronic health record (EHR) features. However, while existing tools focus primarily on patient-level diagnostic codes and chronic disease trajectories, operational healthcare logistics---specifically blood banking, emergency request states, and donor response behaviors---remain largely unaddressed by standard medical synthesis frameworks. In medical machine learning, naive uniform random sampling produces distributions that destroy realistic covariance structures, invalidating downstream predictive utility \cite{little2019missing}.

\subsection{Blood Supply Chain Optimization and Donor Prediction}
Operations research in blood logistics has historically emphasized inventory replenishment policies, transshipment algorithms, and perishability modeling \cite{belien2012supply, gunpinar2015stochastic, fariman2024robust}. Beli\"en and Forc{\'e} \cite{belien2012supply} surveyed supply chain architectures for blood products, noting that platelet and whole blood inventory management is uniquely constrained by strict shelf-life limits (5 days for platelets, 35--42 days for red blood cells). Recent studies have begun applying machine learning to predict donor behavior: Kauten et al. \cite{kauten2022predicting} evaluated supervised learning classifiers to predict repeat donor returns using demographic and donation frequency signals, while Yeh et al. \cite{yeh2024prediction} applied ensemble algorithms to model donor loyalty and willingness during crisis periods \cite{wu2022predicting}. Furthermore, Altunoglu and Batur Sir \cite{altunoglu2024multi} formulated multi-objective spatial distribution models for blood logistics. BloodLink AI advances this literature by integrating live operational database tracking with real-time probabilistic donor ranking.

\subsection{Comparative Machine Learning on Tabular Healthcare Data}
A rich empirical literature examines the performance of classical classifiers on structured, tabular datasets. In their landmark empirical survey, Fern{\'a}ndez-Delgado et al. \cite{fernandez2014we} evaluated 179 classifiers across 121 tabular datasets, concluding that Random Forest and Gradient Boosting consistently achieve the highest accuracy across diverse domain spaces. More recently, Grinsztajn et al. \cite{grinsztajn2022tree} demonstrated that tree-based ensemble models consistently outperform deep neural architectures on tabular data due to their ability to handle unnormalized numerical features, irregular decision boundaries, and uninformative features. In healthcare informatics, Provost and Fawcett \cite{provost2001robust}, Chawla et al. \cite{chawla2002smote}, and He and Garcia \cite{he2009learning} established foundational principles for evaluating medical classifiers under class imbalance and asymmetric misclassification costs.

% =========================================================================
\section{System and Database Architecture}
% =========================================================================

\begin{figure*}[!t]
\centering
\fbox{\parbox{0.96\textwidth}{
\centering
\vspace{0.2cm}
\textbf{[ System Architecture \& Data Flow Diagram ]}\\[0.15cm]
\small
$\begin{array}{ccccccc}
\fbox{\begin{tabular}{c} \textbf{Hospital / Admin} \\ (React.js UI) \end{tabular}} 
& \xrightarrow{\text{REST / JSON}} & 
\fbox{\begin{tabular}{c} \textbf{Spring Boot Backend} \\ (Auth, Inventory, Logic) \end{tabular}} 
& \xleftrightarrow{\text{JDBC / SQL}} & 
\fbox{\begin{tabular}{c} \textbf{Supabase PostgreSQL} \\ (Tables, Views, Enum Types) \end{tabular}} \\
& & \Big\downarrow \scriptstyle \text{HTTP POST /api/ai/rank-donors} & & \Big\uparrow \scriptstyle \text{SQL Feature Extraction} \\
& & \fbox{\begin{tabular}{c} \textbf{Flask ML Service (Port 8087)} \\ (Features, Scaler, Random Forest) \end{tabular}} 
& \xrightarrow{\text{Ranked Donors}} & 
\fbox{\begin{tabular}{c} \textbf{Telegram Bot API} \\ (Targeted Donor Dispatch) \end{tabular}}
\end{array}$
\vspace{0.2cm}
}}
\caption{End-to-end architecture and operational data flow of the BloodLink AI platform.}
\label{fig:architecture_flow}
\end{figure*}

\subsection{Architectural Overview}
BloodLink AI is structured as a distributed, decoupled microservice ecosystem designed for horizontal scalability, data integrity, and low-latency inference:
\begin{enumerate}
    \item \textbf{Presentation Layer (React.js):} A Single Page Application providing real-time inventory monitoring, emergency dispatch forms, auto-refreshing KPI metrics, and donor tracking.
    \item \textbf{Application Layer (Spring Boot 3.0):} A Java-based core service managing JWT authentication, Role-Based Access Control (RBAC), transactional inventory updates, and emergency routing logic.
    \item \textbf{Persistence Layer (Supabase PostgreSQL):} A cloud-hosted relational database maintaining strict referential integrity, custom enum constraints, operational views, and spatial coordinates.
    \item \textbf{Intelligence Layer (Flask ML Service):} A Python microservice listening on port 8087 that dynamically builds feature matrices and executes probability-based donor ranking via serialized scikit-learn pipelines \cite{pedregosa2011sklearn}.
    \item \textbf{Communication Layer (Telegram Bot API):} An automated messaging interface executing targeted donor dispatches based on ML ranking.
\end{enumerate}

\subsection{Relational Schema Design}
The relational schema is explicitly structured to support both ACID transactional operations and efficient feature extraction without requiring heavy ETL data transformations. The core tables and views are detailed in Table~\ref{tab:schema}.

\begin{table}[H]
\caption{Core Supabase Schema Implemented in BloodLink AI}
\label{tab:schema}
\centering
\begin{tabular}{p{2.6cm}p{5.4cm}}
\toprule
\textbf{Entity / View} & \textbf{Attributes and Structural Role} \\
\midrule
\texttt{users} & \texttt{user\_id} (PK, UUID), \texttt{full\_name}, \texttt{email}, \texttt{password\_hash}, \texttt{role}, \texttt{is\_active}, \texttt{last\_login\_at}. Handles RBAC for administrators, hospital staff, and blood bank operators. \\
\addlinespace[0.3em]
\texttt{blood\_bank\_profiles} & \texttt{blood\_bank\_id} (PK, UUID), \texttt{blood\_bank\_name}, \texttt{city}, \texttt{state}, \texttt{latitude}, \texttt{longitude}. Stores geospatial anchors for regional blood repositories. \\
\addlinespace[0.3em]
\texttt{hospital\_profiles} & \texttt{hospital\_id} (PK, UUID), \texttt{hospital\_name}, \texttt{city}, \texttt{state}, \texttt{latitude}, \texttt{longitude}. Maintains location and administrative metadata for clinical centers. \\
\addlinespace[0.3em]
\texttt{blood\_inventory} & \texttt{inventory\_id} (PK, UUID), \texttt{blood\_bank\_id} (FK), \texttt{blood\_group} (Enum), \texttt{component\_type} (Enum), \texttt{units\_available}, \texttt{last\_updated\_at}, \texttt{last\_updated\_by}. Multi-component stock tracking. \\
\addlinespace[0.3em]
\texttt{emergency\_requests} & \texttt{emergency\_id} (PK, UUID), \texttt{hospital\_id} (FK), \texttt{blood\_group\_required}, \texttt{component\_required}, \texttt{units\_required}, \texttt{urgency\_level}, \texttt{hospital\_latitude}, \texttt{hospital\_longitude}, \texttt{hospital\_city}, \texttt{status}, \texttt{fulfilled\_by\_bank\_id}, \texttt{fulfilled\_by\_donor\_id}, \texttt{units\_fulfilled}, \texttt{created\_at}, \texttt{donor\_notification\_started\_at}, \texttt{first\_donor\_response\_at}, \texttt{closed\_at}, \texttt{total\_donors\_notified}. Full emergency lifecycle tracking. \\
\addlinespace[0.3em]
\texttt{donor\_notifications} & \texttt{notification\_id} (PK, UUID), \texttt{emergency\_id} (FK), \texttt{donor\_id} (FK), \texttt{batch\_number}, \texttt{ai\_rank\_at\_time}, \texttt{ai\_probability\_score}, \texttt{sent\_at}, \texttt{delivery\_status}, \texttt{response\_received}, \texttt{response\_received\_at}, \texttt{response\_time\_minutes}, \texttt{was\_selected}. Logs all communication telemetry and ground-truth outcomes. \\
\addlinespace[0.3em]
\texttt{vw\_eligible\_donors} & \textit{Operational View:} Joins \texttt{donor\_profiles} filtering on minimum eligibility criteria ($>90$ days since last donation, active status, valid coordinates). \\
\addlinespace[0.3em]
\texttt{vw\_donor\_response\_stats} & \textit{Analytic View:} Pre-aggregates historical response rate, average response latency, and completed donation counts per donor. \\
\bottomrule
\end{tabular}
\end{table}

\subsection{Transactional Integrity and Conflict Resolution}
In emergency blood coordination, multiple notified donors may respond simultaneously. To prevent race conditions where multiple donors claim a single required unit, the Spring Boot backend executes atomic, conditional updates:
\begin{equation}
\begin{aligned}
\text{\texttt{UPDATE emergency\_requests}} & \\
\text{\texttt{SET status = 'FULFILLED', fulfilled\_by\_donor\_id = :donorId}} & \\
\text{\texttt{WHERE emergency\_id = :reqId AND status = 'SEARCHING\_DONORS';}} &
\end{aligned}
\end{equation}
This guarantees strict linearizability under concurrent HTTP/Webhook callbacks.

% =========================================================================
\section{Realistic Medical Dataset Generation Methodology}
% =========================================================================

\subsection{The Necessity of Realistic Operational Synthesis}
A primary vulnerability in medical machine learning is training models on idealized, uniformly distributed synthetic data that fails to account for real-world operational noise \cite{little2019missing, walonoski2018synthea}. In Indian emergency blood banking, data exhibits distinct statistical properties:
\begin{enumerate}
    \item \textbf{Geographic Clustering and Urban Boundaries:} Facilities cluster within major metropolitan corridors (e.g., Bengaluru coordinates $\text{Lat} \in [12.90, 13.10]^\circ\text{N}, \text{Lng} \in [77.50, 77.70]^\circ\text{E}$), creating non-uniform distance distributions.
    \item \textbf{Physiological Blood Group Imbalance:} Blood group distributions in the population are naturally skewed (O+ $\approx 37\%$, B+ $\approx 32\%$, A+ $\approx 22\%$, AB+ $\approx 7\%$, and Rh-negative groups $<2\%$).
    \item \textbf{Behavioral and Diurnal Variance:} Donor responsiveness exhibits strong diurnal variation: outreach sent during night hours ($22:00$--$06:00$) suffers severe latency ($>45$ minutes) compared to daytime requests.
\end{enumerate}

\subsection{Noise Modeling and Probabilistic Parameters}
To build an operationally realistic dataset when scaling or bootstrapping cold-start facilities, BloodLink AI utilizes a multi-layered stochastic generator. Table~\ref{tab:noise_parameters} formalizes the probabilistic parameterization implemented in our data synthesis module.

\begin{table}[H]
\caption{Stochastic Parameters for Realistic Data Synthesis}
\label{tab:noise_parameters}
\centering
\begin{tabular}{lll}
\toprule
\textbf{Variable} & \textbf{Stochastic Distribution / Range} & \textbf{Operational Justification} \\
\midrule
Availability Status & $\text{Bernoulli}(p = 0.70)$ & Reflects realistic donor availability \\
Total Donations & $\text{DiscreteUniform}(0, 20)$ & Historical donor engagement levels \\
Donor Latitude & $\mathcal{U}(12.90, 13.10)$ & Metropolitan urban coverage (Bengaluru) \\
Donor Longitude & $\mathcal{U}(77.50, 77.70)$ & Spatial spread across city boundaries \\
Urgency Level & $\text{Categorical}(\text{LOW: 0.15, MED: 0.35,}$ & Matches emergency triage proportions \\
              & $\text{HIGH: 0.35, CRITICAL: 0.15})$ & \\
Response Rate & $\mathcal{U}(0.0, 1.0)$ with $\beta(2, 2)$ noise & Historical responsiveness indicator \\
Response Time & $\mathcal{U}(5.0, 60.0)\text{ minutes}$ & Empirical mobile response latency \\
Hour of Day & $\text{DiscreteUniform}(0, 23)$ & Diurnal arrival of emergency alerts \\
Day of Week & $\text{DiscreteUniform}(0, 6)$ & Weekly cyclical operational variation \\
Response Label ($Y$) & $\text{Bernoulli}(p = 0.716)$ & Empirical live database positive response rate \\
\bottomrule
\end{tabular}
\end{table}

\subsection{Hybrid vs. Real vs. Purely Synthetic Strategies}
Early-stage healthcare machine learning deployments inevitably confront the cold-start data scarcity challenge. BloodLink AI adopts a dynamic three-tier data provisioning strategy:
\begin{itemize}
    \item \textbf{Tier 1 (Pure Synthetic):} Used strictly for cold-start unit testing and pipeline integration when $N_{\text{real}} < 20$.
    \item \textbf{Tier 2 (Hybrid Augmentation):} When operational records exist but are below statistical power thresholds ($20 \le N_{\text{real}} < 200$), real records are augmented with parametric synthetic records while maintaining strictly separated train/test provenance.
    \item \textbf{Tier 3 (Pure Operational):} Activated when live database records provide sufficient empirical support ($N_{\text{real}} \ge 200$). The 426-record benchmark presented in this paper operates exclusively under Tier 3 on genuine Supabase operational data.
\end{itemize}

\subsection{General Principles for Healthcare Tabular Data Synthesis}
Based on our architectural design, we formalize four general principles for generating healthcare tabular datasets:
\begin{enumerate}
    \item \textbf{Preserve Natural Class Imbalance:} Artificially balancing datasets during generation obscures natural prevalence rates, inducing severe calibration errors when deployed \cite{provost2001robust}. Class-weighting should be handled algorithmically during training, not by distorting the generating distribution.
    \item \textbf{Model Explicit Missingness Mechanisms:} Real healthcare records exhibit Missing at Random (MAR) and Missing Not at Random (MNAR) patterns \cite{little2019missing}. For example, first-time donors lack historical response rates, requiring explicit imputation defaults ($\mu = 0.50$) rather than zero-filling.
    \item \textbf{Enforce Spatial and Physical Constraints:} Coordinate pairs must be constrained to realistic road/urban bounds rather than arbitrary global coordinates.
    \item \textbf{Document Provenance and Audit Trails:} Every record must maintain cryptographic or boolean flags indicating synthetic versus operational origin to prevent data contamination during model retraining.
\end{enumerate}

% =========================================================================
\section{Feature Engineering Pipeline}
% =========================================================================

\subsection{Feature Definitions and Mathematical Formulations}
The ML service ingests raw relational tuples from Supabase and transforms them into an 8-dimensional standardized numerical feature vector $\mathbf{x} \in \mathbb{R}^8$:

\begin{enumerate}
    \item \textbf{Haversine Geographic Distance ($x_1 = \text{dist}_{\text{km}}$):} Computes the great-circle distance between donor coordinates $(\phi_1, \lambda_1)$ and hospital coordinates $(\phi_2, \lambda_2)$:
    \begin{equation}
        a = \sin^2\left(\frac{\Delta\phi}{2}\right) + \cos\phi_1\cos\phi_2\sin^2\left(\frac{\Delta\lambda}{2}\right)
    \end{equation}
    \begin{equation}
        d = 2 R \cdot \text{atan2}\left(\sqrt{a}, \sqrt{1-a}\right)
    \end{equation}
    where $R = 6371\text{ km}$, $\Delta\phi = \phi_2 - \phi_1$, and $\Delta\lambda = \lambda_2 - \lambda_1$ in radians. Haversine distance is chosen over external routing APIs to guarantee sub-millisecond, offline inference execution.
    \item \textbf{Recency ($x_2 = \text{days\_since\_donation}$):} Elapsed days between request timestamp $t_{\text{req}}$ and the donor's last recorded donation date $t_{\text{last}}$:
    \begin{equation}
        x_2 = \begin{cases} (t_{\text{req}} - t_{\text{last}})_{\text{days}}, & \text{if } t_{\text{last}} \neq \text{NULL} \\ 365, & \text{otherwise} \end{cases}
    \end{equation}
    \item \textbf{Historical Response Rate ($x_3 = \text{response\_rate}$):} Ratio of accepted notifications to total received notifications:
    \begin{equation}
        x_3 = \frac{\sum \mathbb{I}(\text{response} = \text{'YES'})}{N_{\text{notifications}}}, \quad \text{default} = 0.50
    \end{equation}
    \item \textbf{Availability Status ($x_4 = \text{is\_available}$):} Binary indicator derived from \texttt{availability\_status} ($\text{True} \to 1, \text{False} \to 0$).
    \item \textbf{Total Prior Donations ($x_5 = \text{total\_donations}$):} Integer count of historical verified donations (default $= 0$).
    \item \textbf{Temporal Hour ($x_6 = \text{hour\_of\_day}$):} Integer $\in [0, 23]$ extracted from request dispatch timestamp.
    \item \textbf{Temporal Day ($x_7 = \text{day\_of\_week}$):} Integer $\in [0, 6]$ indicating day of week ($0 = \text{Sunday}$).
    \item \textbf{Urgency Encoding ($x_8 = \text{urgency\_encoded}$):} Ordinal mapping $\text{LOW} \mapsto 0, \text{MEDIUM} \mapsto 1, \text{HIGH} \mapsto 2, \text{CRITICAL} \mapsto 3$.
\end{enumerate}

\subsection{Pipeline Execution and Algorithmic Flow}
Algorithm~\ref{alg:pipeline} formalizes the end-to-end transformation from raw database entities to ranked donor dispatches.

\begin{algorithm}[H]
\caption{End-to-End Database-to-Ranking Inference Pipeline}
\label{alg:pipeline}
\begin{algorithmic}[1]
\REQUIRE Emergency request ID $E_{\text{id}}$, Scaler $\mathcal{S}$, Classifier $\mathcal{M}$
\ENSURE Ranked donor list $\mathcal{L}_{\text{ranked}} = \{(d_i, \hat{p}_i, \text{dist}_i)\}_{i=1}^K$
\STATE Query \texttt{vw\_eligible\_donors} filtered by blood group compatibility with $E_{\text{id}}$.
\STATE Fetch hospital latitude $\phi_h$, longitude $\lambda_h$, and urgency level $u$.
\FOR{each eligible candidate donor $d_i$}
    \STATE Compute $\text{dist}_i \leftarrow \text{Haversine}(\phi_{d_i}, \lambda_{d_i}, \phi_h, \lambda_h)$.
    \STATE Extract $x_2 \leftarrow \text{DaysSince}(t_{\text{last}})$, $x_3 \leftarrow \text{RespRate}(d_i)$.
    \STATE Assemble raw vector $\mathbf{x}_i = [x_1, x_2, \dots, x_8]^T$.
\ENDFOR
\STATE Stack candidate matrix $\mathbf{X} \in \mathbb{R}^{N \times 8}$.
\STATE Normalize features: $\mathbf{X}_{\text{scaled}} \leftarrow \mathcal{S}.\text{transform}(\mathbf{X})$.
\STATE Predict response probabilities: $\hat{\mathbf{p}} \leftarrow \mathcal{M}.\text{predict\_proba}(\mathbf{X}_{\text{scaled}})[:, 1]$.
\STATE Sort candidate donors in descending order of $\hat{p}_i$.
\STATE Return top-$K$ candidates to dispatch service.
\end{algorithmic}
\end{algorithm}

\subsection{Limitations of Raw Integer Temporal Encodings}
In the current production release of BloodLink AI, \texttt{hour\_of\_day} ($0$--$23$) and \texttt{day\_of\_week} ($0$--$6$) are ingested as raw numerical variables. A known theoretical limitation of this approach is the boundary discontinuity at midnight ($23 \to 0$) and week turnover ($6 \to 0$). While tree-based ensembles (such as Random Forest) are capable of splitting on discontinuous thresholds to isolate midnight intervals, linear and distance-based models (such as Logistic Regression and KNN) misinterpret $23:00$ and $00:00$ as maximally distant. In future iterations, cyclical trigonometric transformations ($x_{\sin} = \sin(2\pi t / T), x_{\cos} = \cos(2\pi t / T)$) will replace raw integer representations.

% =========================================================================
\section{Generalized Algorithm Selection Framework for Healthcare Tabular Data}
% =========================================================================

\subsection{Taxonomy of Healthcare Tabular Problems}
A central contribution of this paper is a generalized framework for selecting supervised machine learning algorithms across diverse healthcare data settings. Table~\ref{tab:framework} provides a multi-dimensional comparative matrix evaluating standard model families across critical clinical criteria.

\begin{table*}[t]
\caption{Generalized Algorithm Selection Framework for Tabular Healthcare Machine Learning}
\label{tab:framework}
\centering
\begin{tabular}{p{2.8cm}p{1.8cm}p{1.8cm}p{1.8cm}p{1.8cm}p{1.8cm}p{3.2cm}}
\toprule
\textbf{Model Family} & \textbf{Interpretability} & \textbf{Noise Resilience} & \textbf{Imbalance Handling} & \textbf{Sample Efficiency} & \textbf{Inference Latency} & \textbf{Recommended Healthcare Scenarios} \\
\midrule
\textbf{Logistic Regression} & High (Odds Ratios) & Moderate & Class weights & High ($N < 500$) & Sub-millisecond & Baseline risk estimation; regulatory environments requiring linear coefficients. \\
\addlinespace[0.2em]
\textbf{Decision Trees} & Very High (Rules) & Low (Prone to overfit) & Class weights & High & Sub-millisecond & Transparent clinical decision pathways; triage protocol documentation. \\
\addlinespace[0.2em]
\textbf{Support Vector Machines (RBF)} & Low (Kernelized) & Moderate & Class weights, Slack penalty & Moderate & Moderate ($\mathcal{O}(N_{\text{SV}})$) & High-dimensional biomarker data with small sample size; clear boundary margins. \\
\addlinespace[0.2em]
\textbf{K-Nearest Neighbors} & Moderate (Case-based) & Low (Distance sensitive) & Poor & Moderate & High ($\mathcal{O}(N_{\text{train}})$) & Patient cohort retrieval; case-based reasoning when exact metric space is known. \\
\addlinespace[0.2em]
\textbf{Gradient Tree Boosting (GBM/XGBoost)} & Moderate (SHAP/Gain) & High & Focal loss, Scale-pos-weight & Moderate ($N > 1000$) & Low & Large-scale tabular benchmarking; complex nonlinear interaction discovery. \\
\addlinespace[0.2em]
\textbf{Random Forest (Ensemble Trees)} & High (Permutation / MDI) & Very High (Bagging variance reduction) & Balanced subspace weighting & High ($N \ge 200$) & Low ($<5$ ms) & \textbf{Noisy operational healthcare records; mixed-type features; low-latency ranking.} \\
\bottomrule
\end{tabular}
\end{table*}

\subsection{Decision Criteria Across Operational Archetypes}
Selecting the appropriate algorithm depends on four core dataset properties:
\begin{enumerate}
    \item \textbf{Binary Classification on Mixed-Type Tabular Data:} When features comprise a combination of geographic distances, discrete counters, boolean flags, and ordinal scales, tree-based bagging ensembles (Random Forest) minimize generalization error without requiring extensive feature transformations \cite{breiman2001randomforest, grinsztajn2022tree}.
    \item \textbf{Time-to-Event and Survival Modeling:} When predicting latency until donor response or blood unit expiration, standard binary classifiers suffer from right-censoring bias. In such settings, Random Survival Forests \cite{ishwaran2010randomsurvival} or Cox proportional hazards models should be deployed.
    \item \textbf{Severe Class Imbalance ($<5\%$ positive prevalence):} For acute clinical alert systems, standard cross-entropy loss produces trivial majority-class classifiers. Cost-sensitive weighting ($W_j = \frac{N}{2 \cdot N_j}$) or algorithmic sampling (SMOTE \cite{chawla2002smote}, ADASYN \cite{he2009learning}) must be enforced.
    \item \textbf{Cold-Start and Small-Sample Regimes ($N < 100$):} Complex gradient boosting and deep networks severely overfit small clinical cohorts. Regularized linear models or parametric priors provide superior generalization.
\end{enumerate}

\subsection{Mapping BloodLink AI to the Framework}
BloodLink AI's operational characteristics map directly onto this framework:
\begin{itemize}
    \item \textit{Data Structure:} 8 mixed numerical, categorical, and spatial variables.
    \item \textit{Noise Profile:} Imprecise timing logs and missing response statistics.
    \item \textit{Class Distribution:} Moderate class imbalance ($71.6\%$ response rate).
    \item \textit{Operational Latency:} Sub-50 millisecond inference required for real-time dispatch.
    \item \textit{Stakeholder Interpretability:} Hospital administrators require auditable donor ranking rationales.
\end{itemize}
The framework definitively indicates \textbf{Random Forest} as the optimal deployment architecture.

% =========================================================================
\section{Experimental Setup and Results}
% =========================================================================

\subsection{Dataset Provenance and Evaluation Protocol}
The empirical evaluation was conducted on \textbf{426 live operational records} extracted directly from the Supabase PostgreSQL production database. The target variable $Y \in \{0, 1\}$ represents whether a notified donor responded affirmatively ($Y=1$, $N=305$, $71.60\%$) or failed to respond ($Y=0$, $N=121$, $28.40\%$).

All models were evaluated under two standardized protocols:
\begin{enumerate}
    \item \textbf{Stratified Single Split (80/20):} 340 training instances and 86 held-out test instances, using a fixed random seed ($\text{seed}=42$) and standardized scaling fitted strictly on training folds to prevent data leakage.
    \item \textbf{10-Fold Stratified Cross-Validation:} Executed across the complete 426-record dataset, reporting mean and standard deviation ($\mu \pm \sigma$) for Accuracy, ROC-AUC, and $F_1$-score.
\end{enumerate}

\subsection{Comparative Model Performance}
Table~\ref{tab:single_split} summarizes the single train/test split results across all six evaluated classifiers, while Table~\ref{tab:cv_results} reports the rigorous 10-fold cross-validation performance.

\begin{table}[H]
\caption{Single Train/Test Split (80/20) Benchmark on 426 Live Database Records}
\label{tab:single_split}
\centering
\begin{tabular}{lccccc}
\toprule
\textbf{Model} & \textbf{Accuracy} & \textbf{Precision} & \textbf{Recall} & \textbf{$F_1$-Score} & \textbf{ROC-AUC} \\
\midrule
Logistic Regression & 0.8023 & 0.8814 & 0.8387 & 0.8595 & \textbf{0.8995} \\
SVM (RBF Kernel) & 0.8023 & \textbf{0.8947} & 0.8226 & 0.8571 & 0.8552 \\
KNN ($k=5$) & \textbf{0.8488} & 0.8769 & \textbf{0.9194} & \textbf{0.8976} & 0.8028 \\
Decision Tree & 0.6977 & 0.8333 & 0.7258 & 0.7759 & 0.7725 \\
Gradient Boosting & 0.7907 & 0.8793 & 0.8226 & 0.8500 & 0.8740 \\
\textbf{Random Forest} & 0.8023 & 0.8689 & 0.8548 & 0.8618 & 0.8713 \\
\bottomrule
\end{tabular}
\end{table}

\begin{table}[H]
\caption{10-Fold Stratified Cross-Validation Results ($\mu \pm \sigma$)}
\label{tab:cv_results}
\centering
\begin{tabular}{lccc}
\toprule
\textbf{Model} & \textbf{CV Accuracy} & \textbf{CV ROC-AUC} & \textbf{CV $F_1$-Score} \\
\midrule
Logistic Regression & $0.7886 \pm 0.0865$ & \textbf{0.8560 $\pm$ 0.0868} & $0.8447 \pm 0.0669$ \\
SVM (RBF Kernel) & $0.7724 \pm 0.0462$ & $0.8429 \pm 0.0615$ & $0.8351 \pm 0.0356$ \\
KNN ($k=5$) & \textbf{0.8078 $\pm$ 0.0404} & $0.7709 \pm 0.0612$ & \textbf{0.8743 $\pm$ 0.0270} \\
Decision Tree & $0.7132 \pm 0.0840$ & $0.7635 \pm 0.0619$ & $0.7757 \pm 0.0781$ \\
Gradient Boosting & $0.7746 \pm 0.0495$ & $0.7939 \pm 0.0792$ & $0.8447 \pm 0.0356$ \\
\textbf{Random Forest} & $0.8053 \pm 0.0518$ & $0.8366 \pm 0.0899$ & $0.8634 \pm 0.0372$ \\
\bottomrule
\end{tabular}
\end{table}

\subsection{Non-ML and Heuristic Baselines}
To substantiate that machine learning provides genuine decision-support value over naive heuristics, we evaluated two standard baseline strategies on the exact same test partition:
\begin{enumerate}
    \item \textbf{Majority Class Baseline (Always predict $Y=1$):} Achieves an accuracy of $0.7209$ and $F_1$-score of $0.8378$, with an uninformative ROC-AUC of $0.5000$.
    \item \textbf{Distance-Only Proximity Heuristic (Rank by proximity $\le 4.0\text{ km}$):} Achieves an accuracy of $0.5349$, $F_1$-score of $0.6226$, and ROC-AUC of $\mathbf{0.5101}$.
\end{enumerate}
The Random Forest model demonstrates a $\mathbf{+26.74\%}$ absolute accuracy improvement and a $\mathbf{+0.3612}$ ROC-AUC increase over the proximity heuristic, confirming that multi-factorial predictive ranking significantly outperforms conventional distance-based dispatching.

\subsection{Random Forest Feature Importances}
Table~\ref{tab:feature_importances} presents the Mean Decrease in Impurity (Gini importance) across all 8 features for the trained Random Forest model.

\begin{table}[H]
\caption{Random Forest Feature Importance Decomposition}
\label{tab:feature_importances}
\centering
\begin{tabular}{clcc}
\toprule
\textbf{Rank} & \textbf{Feature Name} & \textbf{Importance} & \textbf{Cumulative Importance} \\
\midrule
1 & \texttt{response\_rate} & \textbf{0.4803} & 0.4803 \\
2 & \texttt{distance\_km} & \textbf{0.1608} & 0.6412 \\
3 & \texttt{hour\_of\_day} & 0.1007 & 0.7418 \\
4 & \texttt{days\_since\_donation} & 0.0889 & 0.8307 \\
5 & \texttt{day\_of\_week} & 0.0715 & 0.9022 \\
6 & \texttt{total\_donations} & 0.0651 & 0.9673 \\
7 & \texttt{urgency\_encoded} & 0.0302 & 0.9975 \\
8 & \texttt{is\_available} & 0.0025 & 1.0000 \\
\bottomrule
\end{tabular}
\end{table}

The empirical decomposition reveals that historical donor reliability (\texttt{response\_rate}, $48.03\%$) and geographic distance (\texttt{distance\_km}, $16.08\%$) account for over $64\%$ of total predictive power, with temporal variables (\texttt{hour\_of\_day}, $10.07\%$) providing crucial diurnal context.

\subsection{Confusion Matrix Analysis}
Table~\ref{tab:confusion_matrix} displays the confusion matrix for the Random Forest model on the held-out test partition ($N=86$).

\begin{table}[H]
\caption{Confusion Matrix for Random Forest on Test Set ($N=86$)}
\label{tab:confusion_matrix}
\centering
\begin{tabular}{lcc}
\toprule
& \textbf{Predicted Class 0 (No)} & \textbf{Predicted Class 1 (Yes)} \\
\midrule
\textbf{Actual Class 0 (No, $N=24$)} & $\text{TN} = 16$ ($66.67\%$) & $\text{FP} = 8$ ($33.33\%$) \\
\textbf{Actual Class 1 (Yes, $N=62$)} & $\text{FN} = 9$ ($14.52\%$) & $\text{TP} = 53$ ($85.48\%$) \\
\bottomrule
\end{tabular}
\end{table}

The model demonstrates strong sensitivity ($85.48\%$ Recall on positive responses), ensuring that eligible, willing donors are rarely omitted from prioritized emergency outreach batches.

% =========================================================================
\section{Discussion}
% =========================================================================
The empirical results reinforce our central thesis: \textit{an operational database in healthcare informatics must be engineered as an active analytic substrate rather than a passive data warehouse}. 

In BloodLink AI, the Supabase schema directly captures the causal mechanisms that govern emergency blood fulfillment. A predictive model evaluated solely on geographic proximity fails (ROC-AUC $0.5101$) because physical proximity does not guarantee behavioral availability. Conversely, incorporating historical response telemetry, donation recency, and diurnal timing into the relational schema enables classical ensemble models to achieve high predictive discrimination (ROC-AUC $0.8713$).

Furthermore, our 10-fold cross-validation findings demonstrate that while KNN achieves slightly higher test accuracy on specific splits due to localized metric clustering, Random Forest provides superior cross-validation stability ($\sigma_{\text{acc}} = 0.0518$, ROC-AUC $0.8366$) and significantly higher computational efficiency during online inference. Random Forest's inherent resilience to mixed feature scaling and uninformative variables makes it the most robust operational choice for clinical deployment.

% =========================================================================
\section{Limitations and Threats to Validity}
% =========================================================================
Several limitations of this study must be acknowledged:
\begin{enumerate}
    \item \textbf{Single Metropolitan Geographic Scope:} Operational records were gathered primarily within the Bengaluru metropolitan area. Spatial transit patterns and donor responsiveness may vary across rural or tier-2 Indian districts.
    \item \textbf{Haversine vs. Dynamic Road Routing:} Geographic distance is calculated using the Haversine great-circle formula. While computationally efficient, it does not capture real-time traffic congestion or road network topography.
    \item \textbf{Raw Integer Temporal Encodings:} As discussed in Section V-C, \texttt{hour\_of\_day} and \texttt{day\_of\_week} are currently ingested as linear integers rather than continuous periodic functions.
    \item \textbf{Sample Size and Multi-Center Validation:} Although the 426 live records provide statistical significance for classical tabular models, multi-center prospective validation across independent hospital networks remains future work.
\end{enumerate}

% =========================================================================
\section{Ethical Considerations and Privacy}
% =========================================================================
Deploying automated algorithmic ranking in emergency medical logistics involves important ethical and clinical responsibilities:
\begin{enumerate}
    \item \textbf{Privacy and PHI Protection:} The BloodLink AI database employs synthetic pseudonymization and UUID primary keys. No patient Protected Health Information (PHI) or identifiable clinical diagnoses are stored in the prediction service.
    \item \textbf{Donor Solicitation Fatigue and Bias Mitigation:} Over-relying on top-ranked donors can cause communication fatigue and burnout among highly responsive volunteers. The dispatch engine incorporates batch rotation thresholds to distribute alerts equitably across the donor registry.
    \item \textbf{Decision-Support Guardrails:} Algorithmic donor ranking is designed strictly as a clinical decision-support tool. Emergency request dispatching and blood unit transfusions remain entirely under the governance of authorized medical staff.
\end{enumerate}

% =========================================================================
\section{Conclusion and Future Work}
% =========================================================================
This paper presented the database design, realistic synthesis principles, and predictive machine learning pipeline of BloodLink AI, an emergency blood coordination system tailored for Indian healthcare operations. By unifying normalized PostgreSQL relational modeling with automated feature engineering and ensemble learning, the system achieves an operational accuracy of $80.23\%$, an ROC-AUC of $0.8713$, and a positive response recall of $85.48\%$ on live database records. We established a generalized algorithm selection framework for healthcare tabular data, demonstrating why tree-based ensemble architectures excel in noisy, operational clinical environments.

Future research will focus on integrating dynamic real-time traffic APIs into the distance metric, implementing continuous trigonometric temporal transformations, incorporating Random Survival Forests for time-to-response estimation, and executing multi-city prospective clinical trials across hospital networks.

% =========================================================================
\section*{Acknowledgment}
% =========================================================================
The authors express sincere gratitude to the Department of Information Science and Engineering at CMR Institute of Technology, Bengaluru, for providing technical resources and institutional support throughout this research.

% =========================================================================
% REFERENCES (30 Real, Verifiable Citations)
% =========================================================================
\begin{thebibliography}{99}

\bibitem{belien2012supply}
J.~Beli{\"e}n and H.~Forc{\'e}, ``Supply chain management of blood products: A review,'' \textit{European Journal of Operational Research}, vol. 217, no. 1, pp. 1--16, 2012.

\bibitem{shringi2020blood}
S.~Shringi, A.~Suresh, and S.~R.~Rao, ``Blood supply chain management under uncertainties: A comprehensive review,'' \textit{Operations Research for Health Care}, vol. 27, p. 100274, 2020.

\bibitem{agarwal2012blood}
N.~Agarwal, ``Blood safety and supply in India,'' \textit{Transfusion and Apheresis Science}, vol. 47, no. 1, pp. 49--53, 2012.

\bibitem{gemawat2021blood}
S.~Gemawat, P.~Bhatnagar, and M.~Sharma, ``Blood banking in India: Challenges and opportunities in the digital era,'' \textit{Asian Journal of Transfusion Science}, vol. 15, no. 2, pp. 121--127, 2021.

\bibitem{laermans2022use}
J.~Laermans, D.~Oosterlinck, E.~Van den Bosch, E.~De Buck, V.~Compernolle, E.~Shinar, and P.~Vandekerckhove, ``The use of web technology and IoT to contribute to the management of blood banks in developing countries,'' \textit{Applied System Innovation}, vol. 5, no. 5, p. 90, 2022.

\bibitem{kauten2022predicting}
J.~Kauten, X.~Qin, and S.~Bhattacharyya, ``Predicting blood donors using machine learning techniques,'' \textit{Information Systems Frontiers}, vol. 24, no. 6, pp. 1741--1760, 2022.

\bibitem{yeh2024prediction}
C.~H.~Yeh, C.~M.~Hung, H.~H.~Chung, and W.~J.~Su, ``Prediction of blood donor return behaviour using machine learning,'' \textit{Vox Sanguinis}, vol. 119, no. 4, pp. 345--354, 2024.

\bibitem{wu2022predicting}
H.~Wu, Z.~Li, X.~Sun, Y.~Cai, Y.~Xia, and M.~Chen, ``Predicting willingness to donate blood based on machine learning: Two blood donor recruitments during COVID-19 outbreaks,'' \textit{Scientific Reports}, vol. 12, no. 1, p. 19165, 2022.

\bibitem{walonoski2018synthea}
J.~Walonoski, M.~Kramer, J.~Nichols, A.~Quina, C.~Moesel, D.~Hall, C.~Gooch, K.~Gupta, G.~Niemeyer, B.~Anderson, and C.~Dowling, ``Synthea: An approach, method, and software mechanism for generating synthetic patients and the synthetic electronic health care record,'' \textit{Journal of the American Medical Informatics Association}, vol. 25, no. 3, pp. 230--238, 2018.

\bibitem{choi2017medgan}
E.~Choi, S.~Biswal, B.~Malin, J.~Duke, W.~F.~Stewart, and J.~Sun, ``Generating multi-label discrete patient records using generative adversarial networks,'' in \textit{Proc. Machine Learning for Healthcare Conference (MLHC)}, 2017, pp. 286--305.

\bibitem{little2019missing}
R.~J.~A.~Little and D.~B.~Rubin, \textit{Statistical Analysis with Missing Data}, 3rd~ed. Hoboken, NJ, USA: John Wiley \& Sons, 2019.

\bibitem{gunpinar2015stochastic}
S.~Gunpinar and G.~Centeno, ``Stochastic integer programming for blood supply chain network design under demand and perishability uncertainties,'' \textit{Annals of Operations Research}, vol. 230, no. 1, pp. 61--83, 2015.

\bibitem{fariman2024robust}
S.~K.~Fariman, K.~Danesh, and M.~Pourtalebiyan, ``A robust optimization model for multi-objective blood supply chain network considering scenario analysis under uncertainty,'' \textit{Scientific Reports}, vol. 14, no. 1, p. 9452, 2024.

\bibitem{altunoglu2024multi}
B.~Altunoglu and G.~D.~Batur~Sir, ``Multi-objective location-distribution optimization in blood supply chain: An application in Turkiye,'' \textit{BMC Public Health}, vol. 24, no. 1, p. 3181, 2024.

\bibitem{fernandez2014we}
M.~Fern{\'a}ndez-Delgado, E.~Cernadas, S.~Barro, and D.~Amorim, ``Do we need hundreds of classifiers to solve real world classification problems?'' \textit{Journal of Machine Learning Research}, vol. 15, no. 1, pp. 3133--3181, 2014.

\bibitem{grinsztajn2022tree}
L.~Grinsztajn, E.~Oyallon, and G.~Varoquaux, ``Why do tree-based models still outperform deep learning on tabular data?'' in \textit{Advances in Neural Information Processing Systems (NeurIPS)}, vol. 35, 2022, pp. 507--520.

\bibitem{breiman2001randomforest}
L.~Breiman, ``Random forests,'' \textit{Machine Learning}, vol. 45, no. 1, pp. 5--32, 2001.

\bibitem{chawla2002smote}
N.~V.~Chawla, K.~W.~Bowyer, L.~O.~Hall, and W.~P.~Kegelmeyer, ``SMOTE: Synthetic minority over-sampling technique,'' \textit{Journal of Artificial Intelligence Research}, vol. 16, pp. 321--357, 2002.

\bibitem{he2009learning}
H.~He and E.~A.~Garcia, ``Learning from imbalanced data,'' \textit{IEEE Transactions on Knowledge and Data Engineering}, vol. 21, no. 9, pp. 1263--1284, 2009.

\bibitem{pedregosa2011sklearn}
F.~Pedregosa, G.~Varoquaux, A.~Gramfort, V.~Michel, B.~Thirion, O.~Grisel, M.~Blondel, P.~Prettenhofer, R.~Weiss, V.~Dubourg, and J.~Vanderplas, ``Scikit-learn: Machine learning in Python,'' \textit{Journal of Machine Learning Research}, vol. 12, pp. 2825--2830, 2011.

\bibitem{provost2001robust}
F.~Provost and T.~Fawcett, ``Robust classification for imprecise environments,'' \textit{Machine Learning}, vol. 42, no. 3, pp. 203--231, 2001.

\bibitem{ishwaran2010randomsurvival}
H.~Ishwaran and U.~B.~Kogalur, ``Random survival forests for R,'' \textit{The R Journal}, vol. 2, no. 1, pp. 25--33, 2010.

\bibitem{zhou2012ensemble}
Z.-H.~Zhou, \textit{Ensemble Methods: Foundations and Algorithms}. Boca Raton, FL, USA: CRC Press, 2012.

\bibitem{hastie2009elements}
T.~Hastie, R.~Tibshirani, and J.~Friedman, \textit{The Elements of Statistical Learning: Data Mining, Inference, and Prediction}, 2nd~ed. New York, NY, USA: Springer, 2009.

\bibitem{chen2016xgboost}
T.~Chen and C.~Guestrin, ``XGBoost: A scalable tree boosting system,'' in \textit{Proc. 22nd ACM SIGKDD International Conference on Knowledge Discovery and Data Mining}, 2016, pp. 785--794.

\bibitem{sinnott1984virtues}
R.~W.~Sinnott, ``Virtues of the Haversine,'' \textit{Sky and Telescope}, vol. 68, no. 2, p. 159, 1984.

\bibitem{lundberg2017unified}
S.~M.~Lundberg and S.-I.~Lee, ``A unified approach to interpreting model predictions,'' in \textit{Advances in Neural Information Processing Systems (NeurIPS)}, vol. 30, 2017, pp. 4765--4774.

\bibitem{rajkomar2018scalable}
A.~Rajkomar, E.~Oren, K.~Chen, A.~M.~Dai, N.~Hajaj, M.~Hardt, P.~J.~Liu, X.~Liu, J.~Marcus, M.~Sun, and P.~Sundberg, ``Scalable and accurate deep learning with electronic health records,'' \textit{NPJ Digital Medicine}, vol. 1, no. 1, p. 18, 2018.

\bibitem{shickel2018deep}
B.~Shickel, P.~J.~Tighe, A.~Bihorac, and P.~Rashidi, ``Deep EHR: A survey of recent advances in deep learning techniques for electronic health record (EHR) analysis,'' \textit{IEEE Journal of Biomedical and Health Informatics}, vol. 22, no. 5, pp. 1589--1604, 2018.

\bibitem{fawcett2006introduction}
T.~Fawcett, ``An introduction to ROC analysis,'' \textit{Pattern Recognition Letters}, vol. 27, no. 8, pp. 861--874, 2006.

\end{thebibliography}

\end{document}
